\section{源基准及其适配依据}

本文选择 $\tau$-bench、BFCL、API-Bank 和 AppWorld，是因为它们分别覆盖用户交互、通用函数调用、工具检索与跨应用执行等典型任务形态，且都提供了可追溯的调用证据。但这些证据在原论文中的含义并不相同：有的保存目标动作，有的保存多轮参考调用，有的记录单次 API 请求—结果，有的由最终状态评测器判定任务。本节先回到各基准的任务设计和数据结构，再说明本文从中继承了什么、不能据此声称什么。

\subsection{\texorpdfstring{$\tau$-bench}{tau-bench}}

\subsubsection{任务结构与评测机制}

$\tau$-bench 面向航空和零售客服场景\citep{yao2025taubench}。每个任务由用户指令和一组参考动作（actions）组成：用户指令是一段自然语言，描述用户想完成什么（如取消预订、交换商品、修改地址）；原论文中的 actions 主要包含达到目标状态所需的数据库写操作和需要向用户传达的输出信息，读取操作允许智能体自行决定。零售域有 115 条任务，航空域有 50 条，共 165 条。

任务环境包含结构化数据库：用户档案（email、地址、支付方式、出生日期）、订单、商品、航班和预订记录。智能体在与模拟用户的多轮对话中调用工具，必须遵守退款、改签、身份核验等业务规则。评测不逐字匹配固定轨迹，而是比较执行后的数据库状态是否与标注一致，并结合任务要求检查最终回复是否包含用户所需信息。论文采用 pass$^k$ 衡量重复试验中的稳定成功率。

\subsubsection{适配依据}

本文锁定版本的 \code{tasks\_test.py} 提供了比原论文标注更完整的 typed reference sequence：其中不仅包含写操作，还包含查询（如 \code{get\_user\_details}、\code{get\_reservation\_details}）、搜索（如 \code{search\_direct\_flight}）和计算（如 \code{calculate}）等中间步骤。本文将这些读、写和计算步骤统一作为源参考动作，通过源重放和环境状态检查确认可用性。由于原评测允许智能体通过不同交互过程达到目标状态，这些动作不能自动解释为唯一或最短解，本文不追加最小性声明。

$\tau$-bench 的用户、订单和支付方式都保存在结构化数据库中，用户身份可以由任务标识和数据库记录交叉核对，适合提取用户档案。同时，原生执行揭示出 16 条源参考序列包含业务错误（如查询不存在的商品、对非已交付订单执行换货）。本文没有修补这些序列，而是将其保守过滤。

\subsection{BFCL}

\subsubsection{任务结构与评测机制}

BFCL 是一个综合性函数调用基准，覆盖单轮、并行、多语言和多轮场景\citep{patil2025bfcl}。其多轮部分围绕八组状态化 API（交易、旅行、文件系统、消息、工单等）构造任务，每条任务包含一个多轮查询序列和对应的参考调用轨迹。多轮数据有四个变体：base 是标准多轮任务；miss\_param 故意省略参数以测试追问能力；miss\_func 移除函数以测试澄清能力；long\_context 加长上下文以测试长程准确性。

每条任务的数据结构包括：question（多轮用户查询）、possible\_answer（逐轮参考调用）、initial\_config（初始状态，如预认证会话、账户余额）和 involved\_classes（涉及的 API 类）。评测同时使用 state-based checker 检查最终系统状态是否匹配，以及 response-based checker 检查产生请求结果所必需的调用路径——后者尤其用于无法通过状态变化判断的只读任务。这不等价于对整条参考轨迹进行全局最小性证明。

\subsubsection{适配依据}

本文使用 200 条 base 多轮任务。possible\_answer 保存逐轮参考调用，initial\_config 提供可恢复的初始状态。对于会修改状态的任务，其他调用顺序可能同样正确；对于查询型任务，响应检查会对必要步骤提出更强约束。本文保留原有轮次边界和参考顺序，将其称为"源参考调用"，不把 possible\_answer 解释为经过全量消融的最短 Gold。

BFCL 的主要适配挑战是语义对齐：相同字符串可能在不同 API 类中表示密码、访问令牌、股票代码或普通文本，用户指令中的身份也可能与某个默认 profile 冲突。本文只接受能够由 API 类、参数名、初始状态和指令身份共同支持的字段映射；不能确定时放弃该字段或过滤任务。

\subsection{API-Bank}

\subsubsection{任务结构与评测机制}

API-Bank 将工具调用能力分为三个递进层级\citep{li2023apibank}。Level-1（直接调用）给定 API 定义，评测模型能否按查询正确填参调用，本质是参数填充。Level-2（检索后调用）API 未知且数量大，模型须先用 ToolSearcher 按关键词检索相关 API，再调用单个 API。Level-3（规划—检索—调用）要求模型自主规划并多步调用，难度最高。评测集分别有 214、50 和 50 条对话。

其 73 个 API 由工程师真实实现，涵盖天气、账号管理、日程、健康和财务等领域。对话以多轮形式组织，每个 AI 回合可能产出 API 请求，系统执行后回填结果。每条 API 记录保存请求参数（param\_dict）和返回结果（result），ToolSearcher 是一个特殊 API，输入关键词后返回最相关 API 的元信息。

评测方式为：初始化 API 数据库后，比较预测调用与人工标注调用是否产生相同查询或修改，并检查返回结果是否一致。API 调用使用 Accuracy 评价，调用后的自然语言回复使用 ROUGE-L 评价。这是一种逐调用的一致性检查，而非 AppWorld 那种整段任务的最终状态 evaluator。

\subsubsection{适配依据}

本文使用 214 条 Level-1 对话和 50 条 Level-2 对话，共 264 条源记录。Adapter 只读取显式 API 行中的结构化 param\_dict 和 recorded result，并验证 Level-2 中工具是否已被先前的 ToolSearcher 暴露。Level-3 样例没有与上述两层相同的结构化执行契约，因此没有与其混合。

这里的参考证据粒度是"单次 API 请求—结果"。将同一对话中的多次 API 行连接为一个调用序列，是本文为形成统一评测样本所作的适配，而非 API-Bank 原论文提供的对话级最终状态证明。由此产生的 7 条发布任务可以证明请求记录与上游一致、调用可执行，但不声称具有与 AppWorld 相同强度的最终数据库 oracle。

\subsection{AppWorld}

\subsubsection{任务结构与评测机制}

AppWorld 构建了 9 个日常应用（Gmail、Venmo、Amazon、Spotify、SimpleNote、Splitwise、Todoist、Phone、FileSystem）的沙盒环境，包含 457 个 API、101 张数据库表和约 37 万条数据库记录\citep{trivedi2024appworld}。底层数据由 106 名虚构人物和程序化生成的活动数据填充，模拟虚构人物数字生活的合成数据库。任务让智能体完成跨应用的日常数字任务，如"根据 SimpleNote 里的健身计划播放一个时长够今天锻炼的 Spotify 歌单"。

每条任务包含：任务指令、任务专属数据库副本、期望终态值和评测程序。评测采用基于最终数据库状态的程序化单元测试：计算执行前后的数据库 diff，断言期望的变更全部出现且没有附带损害（如任务没让退货却发起退货）。约 15\% 的任务是问答型，还需核对答案正确性。任务分为 train（105）、dev（60）、test\_normal（168）和 test\_challenge（417）四个 split。

\subsubsection{适配依据}

本文使用 147 条同时具备任务、方案 API 日志（ground\_truth/api\_calls.json）和可执行验证条件的记录，全部来自 dev 和 train split。test\_normal 和 test\_challenge 的任务没有 api\_calls.json，只有 answer.json 和 evaluation.py，不提供标准调用序列。上游 validation solution 的作用是证明任务可解并为数据生成提供参考执行，论文并未把它定义为规范路径或参数最小性证明。

AppWorld 的状态评测能够判断任务目标和数据库副作用，却不直接判断某个已传递参数是否符合当前用途。本文因此另行构造参数级配对调用；其中敏感 optional 参数属于本文的合成扩展，而非 AppWorld 原生接口。

\subsection{证据分层与统一表示}

表\ref{tab:source-characteristics}从任务结构、可见信息、状态环境、参考证据和原始评测五个维度比较四个基准。四个来源最终都会转化为统一调用对象，但统一表示不改变原始证据角色。

\begin{table}[t]
  \centering
  \caption{四个源基准的任务结构与评测机制比较。}
  \label{tab:source-characteristics}
  \scriptsize
  \begin{tabularx}{\textwidth}{L{1.4cm}L{1.6cm}L{1.8cm}L{1.6cm}L{1.8cm}Y}
    \toprule
    来源 & 任务单元 & 智能体可见信息 & 状态环境 & 参考证据 & 原始评测 \\
    \midrule
    $\tau$-bench & 用户模拟任务 & 工具、领域政策、对话 & 航空/零售数据库 & 动作与输出标注 & 数据库状态 + 传达信息 \\
    BFCL & 多轮用户查询 & 函数文档、历史轮次 & 八类状态化 API & \code{possible\_answer} & state + response checker \\
    API-Bank & API 对话 & API 描述或 ToolSearcher & API 数据库与结果 & 人工 API records & 调用一致性 + ROUGE-L \\
    AppWorld & 跨应用数字任务 & ApiDocs、Supervisor、执行历史 & 可重置应用数据库 & validation solution log & DB diff 程序化测试 \\
    \bottomrule
  \end{tabularx}
\end{table}

这种区分直接决定了本文的验证口径。来源重放回答"输出调用是否仍与锁定来源一致"；环境执行回答"调用能否在相应实现中运行"；最终状态验证回答"运行结果是否满足来源任务的状态判据"。三种证据可以同时存在，也可以只具备前两种，任何一种都不自动证明参数已经达到业务意义上的最小集合。
